% Lessons 26--37: arrays, strings, greedy methods, trees, and scheduling.

\patternintoc{Arrays, Strings, and Greedy}
\problemstart{Maximum Subarray}{53}{Medium}{Kadane's algorithm}{Array and running sum}

\subsection{The Problem}
Given an integer array, return the largest sum of a non-empty consecutive
subarray. The chosen values must occupy consecutive positions.

\begin{conceptbox}{Example}
For \texttt{[-2,1,-3,4,-1,2,1,-5,4]}, the best subarray is
\texttt{[4,-1,2,1]}, whose sum is \(6\).
\end{conceptbox}

\subsection{From Repeated Sums to One Scan}
Trying every pair of boundaries repeats almost the same additions. At each
position, only one question matters: is the previous running subarray useful,
or is it better to start again at the current value? A negative prefix can
only reduce every future extension, so it is discarded.

\begin{invariantbox}{the running sum is the best subarray ending here}
After processing index \(i\), \texttt{endingHere} is the maximum sum of a
non-empty subarray whose right endpoint is exactly \(i\). The global answer is
the best value of \texttt{endingHere} seen at any endpoint.
\end{invariantbox}

\begin{visualbox}{Keep a useful prefix or restart}
\centering
\begin{tikzpicture}[font=\small,>={Stealth[length=2.2mm,width=1.55mm]}]
  \foreach \x/\v in {0/-2,1/1,2/-3,3/4,4/-1,5/2,6/1,7/-5,8/4}{
    \node[arraycell] (a\x) at (0.92*\x,0) {\v};
  }
  \node[draw=orange,very thick,fit=(a3)(a4)(a5)(a6),inner sep=1mm] {};
  \draw[decorate,decoration={brace,amplitude=5pt},orange]
    ([yshift=-7mm]a3.south west)--([yshift=-7mm]a6.south east)
    node[midway,below=5mm]{best sum \(=6\)};
  \node[above=7mm of a3,text=teal,font=\bfseries]{restart at 4};
\end{tikzpicture}
\end{visualbox}

\subsection{Pseudocode}
\begin{lstlisting}[style=pseudocode]
endingHere = answer = nums[0]
for each value after the first:
    endingHere = max(value, endingHere + value)
    answer = max(answer, endingHere)
return answer
\end{lstlisting}

\subsection{C++ Implementation}
\begin{lstlisting}[style=compactcode,caption={Complete solution: Maximum Subarray}]
#include <algorithm>
#include <vector>

using namespace std;

class Solution {
public:
    int maxSubArray(vector<int>& nums) {
        int endingHere = nums[0];
        int answer = nums[0];
        for (int i = 1; i < static_cast<int>(nums.size()); ++i) {
            endingHere = max(nums[i], endingHere + nums[i]);
            answer = max(answer, endingHere);
        }
        return answer;
    }
};
\end{lstlisting}

\subsection{Dry Run, Why the Solution Works, and Complexity}
The running sums are \(-2,1,-2,4,3,5,6,1,5\). The maximum is \(6\).
The update rule considers both possible forms of every best subarray ending
at a position: start there or extend the best one ending immediately before.
Time is \(O(n)\), and extra space is \(O(1)\). Start from the first
value rather than zero so arrays containing only negative values work.

\lessonexpansion
  {For \texttt{[-2,1,-3,4,-1,2,1,-5,4]}, the ending sums are
   \(-2,1,-2,4,3,5,6,1,5\). The restart at 4 discards the harmful negative
   prefix, and the global best becomes 6 after processing the following 1.}
  {Every non-empty subarray ending at index \(i\) either begins at \(i\) or
   extends a subarray ending at \(i-1\). Choosing the larger of exactly these
   two possibilities proves the update rule; taking the maximum over all
   endpoints proves the final answer.}
  {Test one negative value, all-negative input, all-positive input, zeros,
   and an optimum that begins at index zero or ends at the final index.}
  {I keep the best sum ending at the current position. A negative accumulated
   prefix is abandoned by comparing extension with a fresh start. The best
   ending value seen anywhere is the answer, in linear time and constant space.}

\chaptercheckpoint
  {A maximum-sum answer must use consecutive array positions.}
  {The best sum ending at the current index and the best overall sum.}
  {Resetting every negative value instead of every negative running prefix.}
  {$O(n)$ time and $O(1)$ extra space.}


\problemstart{Maximum Product Subarray}{152}{Medium}{Dynamic running extremes}{Array}

\subsection{The Problem}
Return the largest product of a non-empty consecutive subarray. Values may be
positive, negative, or zero.

\begin{conceptbox}{Why the minimum matters}
For \texttt{[-2,3,-4]}, the negative product \(-6\) later becomes the best
positive product \(24\) when multiplied by \(-4\).
\end{conceptbox}

\subsection{Building the State}
For sums, one running maximum is enough. Products are different because a
negative multiplier swaps large and small values. So each position
keeps both the maximum and minimum product ending there.

\begin{invariantbox}{both product extremes end at the current position}
After index \(i\), \texttt{high} and \texttt{low} are respectively the
largest and smallest products of subarrays ending exactly at \(i\). Their
three choices are the current value alone, the old high times the value,
and the old low times the value.
\end{invariantbox}

\begin{visualbox}{The old low becomes the new high}
\centering
\begin{tikzpicture}[font=\small,>={Stealth[length=2.2mm,width=1.55mm]}]
  \node[arraycell] (a) {-2}; \node[arraycell,right=0mm of a] (b) {3};
  \node[arraycell,right=0mm of b,active] (c) {-4};
  \node[dsnode,visited,below left=14mm and 24mm of c] (lo)
    {old low times \(-4\)\\\((-6)(-4)=24\)};
  \node[dsnode,below=14mm of c] (hi)
    {old high times \(-4\)\\\((3)(-4)=-12\)};
  \node[dsnode,below right=14mm and 24mm of c] (alone)
    {start here\\\(-4\)};
  \node[below=6mm of lo,text=teal,font=\bfseries]{largest choice: 24};
\end{tikzpicture}
\end{visualbox}

\subsection{Pseudocode and C++}
\begin{lstlisting}[style=pseudocode]
low = high = answer = nums[0]
for value in nums[1...]:
    oldHigh = high; oldLow = low
    high = max(value, oldHigh*value, oldLow*value)
    low  = min(value, oldHigh*value, oldLow*value)
    answer = max(answer, high)
\end{lstlisting}

\begin{lstlisting}[style=compactcode,caption={Complete solution: Maximum Product Subarray}]
#include <algorithm>
#include <vector>

using namespace std;

class Solution {
public:
    int maxProduct(vector<int>& nums) {
        int high = nums[0], low = nums[0], answer = nums[0];
        for (int i = 1; i < static_cast<int>(nums.size()); ++i) {
            const int oldHigh = high, oldLow = low, value = nums[i];
            high = max({value, oldHigh * value, oldLow * value});
            low = min({value, oldHigh * value, oldLow * value});
            answer = max(answer, high);
        }
        return answer;
    }
};
\end{lstlisting}

\subsection{Time, Space, and Edge Cases}
Each value performs constant work, so time is \(O(n)\) and space is \(O(1)\).
Zero automatically starts a new product. Save both old extremes before either
is overwritten. Test one value, all negatives, zeros between useful segments,
and an odd or even count of negative values.

\lessonexpansion
  {For \texttt{[-2,3,-4]}, start with \(high=low=-2\). At value 3, the
   choices are \(3,-6,-6\), so \(high=3\) and \(low=-6\). At value
   \(-4\), the choices are \(-4,-12,24\); the old minimum becomes the new
   maximum, producing the answer \(24\).}
  {Every product ending at the current index either starts at that value or
   extends a product ending at the previous index. Multiplication by a
   negative value can exchange the maximum and minimum, which is why keeping
   both extremes is enough and needed.}
  {Check \texttt{[-2]}, \texttt{[0,2]}, \texttt{[-2,0,-1]},
   \texttt{[-2,3,-4]}, and a sequence containing two separated zeroes.}
  {I track both the largest and smallest product ending at each position.
   The smallest is retained because a later negative number can turn it into
   the largest. Each value has three choices, giving linear time and
   constant extra space.}

\chaptercheckpoint
  {Signs can reverse which running product is useful.}
  {The maximum and minimum product ending at each position.}
  {Updating the minimum from an already modified maximum.}
  {$O(n)$ time and $O(1)$ space.}


\problemstart{Longest Repeating Character Replacement}{424}{Medium}{Sliding window}{Frequency array}

\subsection{The Problem}
Given an uppercase string and an integer \(k\), return the longest substring
that can be made of one repeated character after replacing at most \(k\)
characters.

\subsection{Window Condition}
Inside a window, keep the most frequent character and replace every other
position. A window is valid exactly when \(\text{window length}-\text{maximum frequency}\le k.\)
When this cost becomes too large, move the left boundary.

\begin{invariantbox}{the active window needs at most k replacements}
After shrinking, the window satisfies the replacement budget. The frequency
array describes its characters, while \texttt{maxFrequency} is a safe
nondecreasing upper bound used to decide when the left boundary must move.
\end{invariantbox}

\begin{visualbox}{Change one B to make four A's}
\centering
\begin{tikzpicture}[font=\small,>={Stealth[length=2.2mm,width=1.55mm]}]
  \foreach \x/\v in {0/A,1/A,2/B,3/A,4/B,5/B,6/A}{
    \node[arraycell] (s\x) at (1.0*\x,0) {\v};
  }
  \node[draw=orange,very thick,fit=(s0)(s1)(s2)(s3),inner sep=1mm] {};
  \node[below=9mm of s2,text=orange] (change) {change this B to A};
  \draw[flowarrow,orange] (change.north)--(s2.south);
  \node[right=12mm of s6,align=left,text=navy]
    {four positions, three already A\\so only one change is needed};
  \node[above=7mm of s0,text=teal]{left};
  \node[above=7mm of s3,text=orange]{right};
\end{tikzpicture}
\end{visualbox}

\subsection{Pseudocode and C++}
\begin{lstlisting}[style=pseudocode]
left = 0; maxFrequency = 0; answer = 0
for right from 0 to n-1:
    increment count[s[right]]
    update maxFrequency
    while windowLength - maxFrequency > k:
        decrement count[s[left++]]
    update answer with windowLength
\end{lstlisting}

\begin{lstlisting}[style=compactcode,caption={Complete solution: Character Replacement}]
#include <algorithm>
#include <array>
#include <string>

using namespace std;

class Solution {
public:
    int characterReplacement(string s, int k) {
        array<int, 26> count{};
        int left = 0, maxFrequency = 0, answer = 0;
        for (int right = 0; right < static_cast<int>(s.size()); ++right) {
            maxFrequency = max(maxFrequency, ++count[s[right] - 'A']);
            while (right - left + 1 - maxFrequency > k) {
                --count[s[left] - 'A'];
                ++left;
            }
            answer = max(answer, right - left + 1);
        }
        return answer;
    }
};
\end{lstlisting}

\subsection{Why a Stale Maximum Is Safe}
The stored maximum frequency need not decrease when the left side moves. A
stale larger value may delay shrinking, but it cannot create a new answer
longer than a window that was already achievable when that frequency was
observed. Both boundaries move only forward: \(O(n)\) time and \(O(1)\) space.

\lessonexpansion
  {For \texttt{AABABBA} with \(k=1\), the window \texttt{AABA} has length
   four and maximum frequency three, so it needs one replacement. Extending
   to \texttt{AABAB} needs two replacements; move the left boundary until the
   window again satisfies \(length-maxFrequency\le1\). The best length remains
   four.}
  {For a fixed window, changing every other character is best, so
   the required replacements equal its length minus its largest frequency.
   The algorithm records every maximal valid window reached by the right
   boundary.}
  {Check \texttt{"A"}, \(k=0\); all-identical text; alternating characters;
   \(k\) at least the string length; and a best window ending at the final
   character.}
  {I use a variable sliding window. Character counts identify the dominant
   letter, and I shrink only when the replacement cost exceeds \(k\). Since
   both pointers move forward, the running time is linear.}

\chaptercheckpoint
  {Find the longest substring repairable with at most $k$ changes.}
  {Two boundaries, 26 counts, and the largest frequency seen in the window.}
  {Using the number of distinct characters instead of replacement cost.}
  {$O(n)$ time and $O(1)$ alphabet space.}


\problemstart{Longest Palindromic Substring}{5}{Medium}{Expand around centers}{String boundaries}

\subsection{The Problem}
Return the longest consecutive substring that reads the same from left to
right and right to left.

\subsection{Every Palindrome Has a Center}
There are \(n\) odd-length centers on characters and \(n-1\) even-length
centers between characters. Expand outward while the two characters match;
every palindrome is discovered from its unique center.

\begin{invariantbox}{the characters strictly inside the pointers match}
During an expansion, the substring between \texttt{left+1} and
\texttt{right-1} is already a palindrome. Equal boundary characters extend it
by two; the first mismatch proves that center cannot grow further.
\end{invariantbox}

\begin{visualbox}{A palindrome grows from a letter or a gap}
\centering
\begin{tikzpicture}[font=\small,>={Stealth[length=2.2mm,width=1.55mm]}]
  \foreach \x/\v in {0/b,1/a,2/b,3/a,4/d}{
    \node[arraycell] (p\x) at (1.1*\x,1.0) {\v};
  }
  \node[draw=orange,very thick,fit=(p0)(p1)(p2),inner sep=1mm] {};
  \node[right=12mm of p4,text=orange]{middle letter: \texttt{a}};

  \foreach \x/\v in {0/c,1/b,2/b,3/d}{
    \node[arraycell] (q\x) at (1.1*\x,-0.7) {\v};
  }
  \node[draw=teal,very thick,fit=(q1)(q2),inner sep=1mm] {};
  \node[right=12mm of q3,text=teal]{middle gap: between the two b's};
\end{tikzpicture}
\end{visualbox}

\subsection{Pseudocode}
\begin{lstlisting}[style=pseudocode]
bestStart = 0; bestLength = 1
for center from 0 to n-1:
    expand(center, center)       // odd length
    expand(center, center + 1)   // even length
expand(left, right):
    while boundaries match: move both outward
    update the longest valid range just crossed
return the recorded substring
\end{lstlisting}

\subsection{C++ Implementation}
\begin{lstlisting}[style=compactcode,caption={Complete solution: Longest Palindromic Substring}]
#include <string>

using namespace std;

class Solution {
public:
    string longestPalindrome(string s) {
        int bestStart = 0, bestLength = 1;
        auto expand = [&](int left, int right) {
            while (left >= 0 && right < static_cast<int>(s.size()) &&
                   s[left] == s[right]) {
                --left;
                ++right;
            }
            const int length = right - left - 1;
            if (length > bestLength) {
                bestLength = length;
                bestStart = left + 1;
            }
        };

        for (int center = 0; center < static_cast<int>(s.size()); ++center) {
            expand(center, center);
            expand(center, center + 1);
        }
        return s.substr(bestStart, bestLength);
    }
};
\end{lstlisting}

\subsection{Dry Run and Complexity}
For \texttt{babad}, expanding around index 1 finds \texttt{bab}; expanding
around index 2 finds \texttt{aba}. Either maximum is valid. There are
\(2n-1\) centers and each may expand \(O(n)\), so time is \(O(n^2)\) and
extra space is \(O(1)\).

\lessonexpansion
  {For \texttt{cbbd}, the character centers produce only length-one
   palindromes. The gap between the two \texttt{b} characters expands once to
   form \texttt{bb}; the next outward comparison \texttt{c} versus
   \texttt{d} fails, so \texttt{bb} is recorded.}
  {Every palindrome has one unique center: a character for odd length or a
   gap for even length. Expansion enumerates every palindrome belonging to
   that center, so testing all \(2n-1\) centers cannot miss the optimum.}
  {Check empty handling if the interface permits it, one character, two equal
   characters, two unequal characters, an even optimum, and tied optima.}
  {I expand from every character and every adjacent-character gap. Matching
   boundaries extend a valid palindrome; after the first mismatch I compare
   the last valid length with the best recorded range.}

\chaptercheckpoint
  {The answer is a consecutive palindrome, not a subsequence.}
  {A center and two outward-moving boundaries.}
  {Checking only odd centers and missing even palindromes.}
  {$O(n^2)$ time and $O(1)$ extra space.}


\problemstart{Longest Substring Without Repeating Characters}{3}{Medium}{Sliding window}{Last-position table}

\subsection{The Problem}
Return the length of the longest substring containing no repeated character.

\subsection{Jump the Left Boundary}
When the right character was already seen inside the active window, the next
valid window must begin after its previous occurrence. A last-position table
lets the left boundary jump directly rather than removing characters one by
one.

\begin{invariantbox}{the active window contains unique characters}
Before recording the answer at each right endpoint, \texttt{left} is greater
than the previous position of every duplicated character in the window.
So \texttt{s[left..right]} contains no repeated character.
\end{invariantbox}

\begin{visualbox}{A repeated a moves the left edge past the old a}
\centering
\begin{tikzpicture}[font=\small,>={Stealth[length=2.2mm,width=1.55mm]}]
  \foreach \x/\v in {0/a,1/b,2/c,3/a,4/b,5/c,6/b,7/b}{
    \node[arraycell] (u\x) at (0.95*\x,0) {\v};
  }
  \node[draw=teal,very thick,fit=(u1)(u2)(u3),inner sep=1mm] {};
  \node[below=7mm of u0,text=orange]{old a};
  \node[below=7mm of u3,text=orange]{repeated a};
  \draw[->,very thick,orange] ([yshift=9mm]u0.north)--
    node[edgelabel,above]{move the left edge here} ([yshift=9mm]u1.north);
  \node[above=16mm of u2,text=teal,font=\bfseries]{valid window: \texttt{bca}};
\end{tikzpicture}
\end{visualbox}

\subsection{Pseudocode}
\begin{lstlisting}[style=pseudocode]
last position of every character = -1
left = 0; answer = 0
for right from 0 to n-1:
    left = max(left, last[s[right]] + 1)
    last[s[right]] = right
    answer = max(answer, right - left + 1)
return answer
\end{lstlisting}

\subsection{C++ Implementation}
\begin{lstlisting}[style=compactcode,caption={Complete solution: Longest Substring Without Repeats}]
#include <algorithm>
#include <array>
#include <string>

using namespace std;

class Solution {
public:
    int lengthOfLongestSubstring(string s) {
        array<int, 256> last;
        last.fill(-1);
        int left = 0, answer = 0;
        for (int right = 0; right < static_cast<int>(s.size()); ++right) {
            const unsigned char ch = static_cast<unsigned char>(s[right]);
            left = max(left, last[ch] + 1);
            last[ch] = right;
            answer = max(answer, right - left + 1);
        }
        return answer;
    }
};
\end{lstlisting}

\subsection{Why It Works and Complexity}
The left boundary never moves backward, even if the character's last
occurrence lies before the current window. Each character is processed once,
so time is \(O(n)\) and the fixed table uses \(O(1)\) space for byte-valued
characters. Empty input correctly returns zero.

\lessonexpansion
  {For \texttt{abba}, the second \texttt{b} moves \texttt{left} from 0 to 2,
   leaving \texttt{b}. At the final \texttt{a}, its previous position 0 lies
   outside the current window, so \texttt{left} stays 2 and the valid window
   becomes \texttt{ba}. The earlier window \texttt{ab} remains the best.}
  {The left boundary is always one position beyond the most recent duplicate
   that lies inside the active window. So every recorded window is
   unique, and every longest unique suffix ending at each right boundary is
   considered.}
  {Test an empty string, all identical characters, all unique characters,
   repeated characters outside the current window, and byte values with a
   negative plain-\texttt{char} representation.}
  {I store the latest index of each character. When a repeat appears, the left
   boundary jumps forward past its previous copy but never moves backward.
   Each character is processed once, so the algorithm is linear.}

\chaptercheckpoint
  {A longest substring must contain unique characters.}
  {Two boundaries and the most recent index of each character.}
  {Moving left backward to an occurrence outside the current window.}
  {$O(n)$ time and $O(1)$ table space.}


\problemstart{Palindromic Substrings}{647}{Medium}{Interval dynamic programming}{Two-dimensional DP table}

\subsection{The Problem}
Count every palindromic substring. Equal text at different positions counts
as different substrings.

\begin{conceptbox}{Examples and positional counting}
\texttt{"abc"} contains three palindromic substrings: \texttt{"a"},
\texttt{"b"}, and \texttt{"c"}. The string \texttt{"aaa"} contains six:
three length-one intervals, two length-two intervals, and one length-three
interval. Equal text at different positions counts separately.
\end{conceptbox}

\subsection{From Repeated Checks to Interval DP}
A direct method can generate every substring and scan inward to determine
whether it is a palindrome. There are \(O(n^2)\) intervals and each check can
cost \(O(n)\), producing \(O(n^3)\) work.

Instead, let \(dp[i][j]\) mean that the interval from index \(i\) through
index \(j\) is a palindrome. The outer characters must match, and the interior
must already be a palindrome. Intervals of length one and two form the base
cases: \(dp[i][j]=(s[i]=s[j])\land\bigl(j-i\le1\ \lor\ dp[i+1][j-1]\bigr).\)

\begin{invariantbox}{every interior interval is solved before its exterior}
The table is filled with \(i\) moving from right to left and \(j\) moving from
\(i\) to the right. So \(dp[i+1][j-1]\) is known before
\(dp[i][j]\). Every true table cell corresponds to exactly one pair of
substring endpoints and increases the answer once.
\end{invariantbox}

\begin{visualbox}{Six checked cells mean six palindromes in \texttt{aaa}}
\centering
\begin{tikzpicture}[font=\small,>={Stealth[length=2.2mm,width=1.55mm]}]
  \matrix[matrix of nodes,
    nodes={arraycell,minimum width=13mm,minimum height=10mm},
    row sep=-\pgflinewidth,column sep=-\pgflinewidth] (m) {
      \checkmark & \checkmark & \checkmark \\
                 & \checkmark & \checkmark \\
                 &            & \checkmark \\
  };
  \node[above=3mm of m-1-1]{end 0};
  \node[above=3mm of m-1-2]{end 1};
  \node[above=3mm of m-1-3]{end 2};
  \node[left=3mm of m-1-1]{start 0};
  \node[left=3mm of m-2-2]{start 1};
  \node[left=3mm of m-3-3]{start 2};
  \node[right=18mm of m-2-3,dsnode,fill=softteal]
    {each check mark is\\one palindrome range};
\end{tikzpicture}
\end{visualbox}

\subsection{Pseudocode}
\begin{lstlisting}[style=pseudocode]
count = 0
dp = n by n table filled with false
for left from n - 1 down to 0:
    for right from left to n - 1:
        shortInterior = (right - left <= 1)
        if s[left] == s[right] and
           (shortInterior or dp[left + 1][right - 1]):
            dp[left][right] = true
            count++
return count
\end{lstlisting}

\subsection{C++ Implementation}
\begin{lstlisting}[style=compactcode,caption={Complete solution: Palindromic Substrings}]
#include <string>
#include <vector>

using namespace std;

class Solution {
public:
    int countSubstrings(string s) {
        const int n = static_cast<int>(s.size());
        vector<vector<bool>> dp(
            n, vector<bool>(n, false));
        int answer = 0;

        for (int left = n - 1; left >= 0; --left) {
            for (int right = left; right < n; ++right) {
                const bool shortInterior = right - left <= 1;
                if (s[left] == s[right] &&
                    (shortInterior || dp[left + 1][right - 1])) {
                    dp[left][right] = true;
                    ++answer;
                }
            }
        }
        return answer;
    }
};
\end{lstlisting}

\subsection{Detailed Visual Dry Run}
\begin{dryrunbox}{Fill order for \texttt{aaa}}
\begin{tabularx}{\linewidth}{@{}c c X c@{}}
\toprule
\(i\) & \(j\) & Reason & Count \\
\midrule
2 & 2 & Single character \texttt{a} & 1 \\
1 & 1 & Single character \texttt{a} & 2 \\
1 & 2 & Equal endpoints and length two & 3 \\
0 & 0 & Single character \texttt{a} & 4 \\
0 & 1 & Equal endpoints and length two & 5 \\
0 & 2 & Equal endpoints; \(dp[1][1]\) is true & 6 \\
\bottomrule
\end{tabularx}
\end{dryrunbox}

\subsection{Why It Works, Complexity, and Edge Cases}
The update rule is needed because a palindrome requires matching endpoints
and a palindromic interior. It is enough for the same reason. Since every
interval is represented by one table cell, every palindromic substring is
counted exactly once. Time and table space are both \(O(n^2)\).

Test one character, all equal characters, no multi-character palindrome, and
both odd- and even-length palindromes. The loop order must not be reversed,
because the interior state must already be available.

\lessonexpansion
  {For \texttt{abba}, the single-character diagonal is true first. Length-two
   intervals then identify \texttt{bb}. When the cell for indices \((0,3)\)
   is reached, the endpoints are both \texttt{a} and the interior cell for
   \texttt{bb} is already true, so \texttt{abba} is counted. Each true cell
   represents one distinct pair of positions.}
  {The update rule is needed: unequal endpoints cannot form a palindrome,
   and a long interval with matching endpoints still fails if its interior is
   not palindromic. It is enough because matching endpoints placed around
   a palindromic interior produce a palindrome. The fill order guarantees that
   the required interior result is available.}
  {Check an empty string if the surrounding interface permits it, one
   character, two equal characters, two unequal characters, all-equal input,
   and strings containing both even- and odd-length palindromes.}
  {I assign one Boolean state to every substring interval. Matching endpoints
   form a palindrome when the interval is shorter than three or its interior
   was already proved palindromic. Filling from right to left resolves that
   dependency and counts each positional substring once.}

\chaptercheckpoint
  {Count all palindromic intervals rather than the longest one.}
  {A Boolean table whose cells represent substring endpoints.}
  {Reading an interior cell before it has been computed.}
  {$O(n^2)$ time and $O(n^2)$ space.}


\patternintoc{Interval Scheduling}
\problemstart{Meeting Rooms}{252}{Easy}{Sort and scan}{Intervals}

\subsection{The Problem}
Given meeting intervals, determine whether one person can attend all of them.
Meetings conflict when their occupied times overlap.

\subsection{Sort by Start Time}
After sorting, any overlap must occur between two adjacent meetings: if the
current start is earlier than the previous end, both require the room at the
same time.

\begin{invariantbox}{all meetings before the current one are compatible}
During the scan, every already checked adjacent pair is non-overlapping. The
previous interval has the latest relevant start among them, so comparing its
end with the current start is enough.
\end{invariantbox}

\begin{visualbox}{The first two meetings overlap}
\centering
\begin{tikzpicture}[font=\small,>={Stealth[length=2.2mm,width=1.55mm]},x=7mm]
  \draw[->,navy] (0,0)--(11,0);
  \foreach \x in {0,2,4,6,8,10} \draw (\x,0.1)--(\x,-0.1) node[edgelabel,below]{\x};
  \draw[line width=4pt,teal] (0,0.7)--(3,0.7) node[edgelabel,midway,above]{[0,3]};
  \draw[line width=4pt,orange] (2,1.4)--(6,1.4) node[edgelabel,midway,above]{[2,6]};
  \draw[line width=4pt,navy] (7,0.7)--(10,0.7) node[edgelabel,midway,above]{[7,10]};
  \node[text=orange] at (2.5,2.1) {both use time 2 to 3};
\end{tikzpicture}
\end{visualbox}

\subsection{Pseudocode}
\begin{lstlisting}[style=pseudocode]
sort intervals by start time
for i from 1 to n-1:
    if intervals[i].start < intervals[i-1].end:
        return false
return true
\end{lstlisting}

\subsection{C++ Implementation}
\begin{lstlisting}[style=compactcode,caption={Complete solution: Meeting Rooms}]
#include <algorithm>
#include <vector>

using namespace std;

class Solution {
public:
    bool canAttendMeetings(vector<vector<int>>& intervals) {
        sort(intervals.begin(), intervals.end());
        for (int i = 1; i < static_cast<int>(intervals.size()); ++i) {
            if (intervals[i][0] < intervals[i - 1][1]) return false;
        }
        return true;
    }
};
\end{lstlisting}

\subsection{Complexity and Boundary Rule}
Sorting costs \(O(n\log n)\); the scan is \(O(n)\). Adjacent meetings such as
\([1,2]\) and \([2,3]\) do not overlap, so the comparison is strict
\texttt{<}, not \texttt{<=}.

\lessonexpansion
  {Sort \texttt{[[7,10],[2,4],[1,3]]} into
   \texttt{[[1,3],[2,4],[7,10]]}. The second meeting starts at 2 before the
   previous meeting ends at 3, so the schedule immediately fails.}
  {After sorting by start time, if any overlap exists, one also exists between
   two adjacent intervals in that order. Checking each interval against the
   preceding end is enough.}
  {Check no meetings, one meeting, touching endpoints, identical intervals,
   nested intervals, and input already sorted in reverse order.}
  {I sort by start time so conflicts become adjacent. I then compare each
   start with the preceding end; a strict earlier start means overlap.}

\chaptercheckpoint
  {Determine whether any time intervals overlap.}
  {Intervals sorted by start time and the previous end.}
  {Treating a meeting that starts exactly at the prior end as a conflict.}
  {$O(n\log n)$ time.}


\patternintoc{Greedy Reachability}
\problemstart{Jump Game}{55}{Medium}{Greedy reachability}{Array}

\subsection{The Problem}
Each array value is the maximum jump length from that position. Return whether
the final index is reachable from index zero.

\subsection{Keep the Reachable Frontier}
If index \(i\) is within the farthest reachable position, then every landing
position up to \(i+\texttt{nums[i]}\) also becomes reachable. If the scan ever
passes the frontier, no future position can be used.

\begin{invariantbox}{every index through farthest is reachable}
Before examining index \(i\), every position from zero through
\texttt{farthest} can be reached. When \(i\le\texttt{farthest}\), processing
its jump can only extend this consecutive reachable prefix.
\end{invariantbox}

\begin{visualbox}{Reach index 2, then extend the reach to index 4}
\centering
\begin{tikzpicture}[font=\small,>={Stealth[length=2.2mm,width=1.55mm]}]
  \foreach \x/\v in {0/2,1/3,2/1,3/1,4/4}{
    \node[arraycell] (j\x) at (1.15*\x,0) {\v};
    \node[below=1mm of j\x,font=\scriptsize]{\x};
  }
  \draw[->,very thick,teal] ([yshift=9mm]j0.north)--node[edgelabel,above]{from 0, reach 2} ([yshift=9mm]j2.north);
  \draw[->,very thick,orange] ([yshift=16mm]j1.north)--node[edgelabel,above]{from 1, reach 4} ([yshift=16mm]j4.north);
\end{tikzpicture}
\end{visualbox}

\subsection{Pseudocode}
\begin{lstlisting}[style=pseudocode]
farthest = 0
for i from 0 to n-1:
    if i > farthest: return false
    farthest = max(farthest, i + nums[i])
    if farthest >= n - 1: return true
return true
\end{lstlisting}

\subsection{C++ Implementation}
\begin{lstlisting}[style=compactcode,caption={Complete solution: Jump Game}]
#include <algorithm>
#include <vector>

using namespace std;

class Solution {
public:
    bool canJump(vector<int>& nums) {
        int farthest = 0;
        for (int i = 0; i < static_cast<int>(nums.size()); ++i) {
            if (i > farthest) return false;
            farthest = max(farthest, i + nums[i]);
            if (farthest >= static_cast<int>(nums.size()) - 1) return true;
        }
        return true;
    }
};
\end{lstlisting}

\subsection{Complexity and Mistakes}
The scan is \(O(n)\) time and \(O(1)\) space. The state is reachability, not
the exact path or the minimum number of jumps. A zero is harmless if the
frontier already extends beyond it, but fatal when it blocks the frontier.

\lessonexpansion
  {For \texttt{[3,2,1,0,4]}, indices 0 through 3 are reachable, but none can
   extend the frontier beyond 3. When the scan reaches index 4, \(4>3\), so
   the destination is unreachable. For \texttt{[2,3,1,1,4]}, index 1 extends
   the frontier directly to 4.}
  {The reachable indices always form a prefix. Processing every index inside
   that prefix records the farthest next endpoint; seeing an index
   beyond it proves no earlier jump can reach that position or anything after
   it.}
  {Test one element, a leading zero, zeros already covered by the frontier,
   an exact jump to the end, and a frontier that stops one position short.}
  {I do not choose an exact path. I scan the reachable prefix and keep its
   farthest extension. If the scan ever moves beyond that frontier the answer
   is false; reaching the last index makes it true.}

\chaptercheckpoint
  {Each position supplies a maximum forward reach.}
  {The farthest reachable index.}
  {Assuming every zero is automatically impossible.}
  {$O(n)$ time and $O(1)$ space.}


\patternintoc{BSTs and Tree Depth}
\problemstart{Maximum Depth of Binary Tree}{104}{Easy}{Recursive DFS}{Binary tree}

\subsection{The Problem}
Return the number of nodes on the longest path from the root to a leaf. An
empty tree has depth zero.

\subsection{Recursive Definition}
The maximum depth of a non-null node is one plus the larger depth of its two
subtrees. This definition translates directly into postorder recursion.

\begin{invariantbox}{each call returns the exact depth of its subtree}
After both child calls return, their depths are correct by the recursive
hypothesis. Adding one for the current node produces the correct
longest root-to-leaf node count for the current subtree.
\end{invariantbox}

\begin{visualbox}{The deepest root-to-leaf path contains three nodes}
\centering
\begin{tikzpicture}[font=\small,>={Stealth[length=2.2mm,width=1.55mm]}]
  \node[trienode,active] (r) {3};
  \node[trienode,below left=12mm and 18mm of r] (a) {9};
  \node[trienode,below right=12mm and 18mm of r] (b) {20};
  \node[trienode,visited,below left=12mm and 8mm of b] (c) {15};
  \node[trienode,below right=12mm and 8mm of b] (d) {7};
  \draw[-,thick,navy] (r)--(a); \draw[-,very thick,orange] (r)--(b);
  \draw[-,very thick,orange] (b)--(c); \draw[-,thick,navy] (b)--(d);
  \node[right=25mm of r,align=left] (steps)
    {15 has depth 1\\20 has depth 2\\3 has depth 3};
  \node[below=6mm of c,text=orange,font=\bfseries]{deepest path};
\end{tikzpicture}
\end{visualbox}

\subsection{Pseudocode}
\begin{lstlisting}[style=pseudocode]
depth(node):
    if node is null: return 0
    leftDepth = depth(node.left)
    rightDepth = depth(node.right)
    return 1 + max(leftDepth, rightDepth)
\end{lstlisting}

\subsection{C++ Implementation}
\begin{lstlisting}[style=compactcode,caption={Complete solution: Maximum Depth of Binary Tree}]
#include <algorithm>

using namespace std;

class Solution {
public:
    int maxDepth(TreeNode* root) {
        if (root == nullptr) return 0;
        return 1 + max(maxDepth(root->left), maxDepth(root->right));
    }
};
\end{lstlisting}

\subsection{Time, Space, and Edge Cases}
Every node is visited once: \(O(n)\) time. The recursion stack is \(O(h)\),
which is \(O(n)\) for a skewed tree and \(O(\log n)\) for a balanced tree.

\lessonexpansion
  {In the tree \texttt{[3,9,20,null,null,15,7]}, both children of 20 return
   depth 1, so node 20 returns 2. Node 9 returns 1, and the root returns
   \(1+\max(1,2)=3\).}
  {The empty subtree has depth zero. Assuming both child calls return correct
   depths, adding one to their maximum gives exactly the longest root-to-leaf
   node count, which proves the update rule by working upward from the leaves.}
  {Check a null root, one node, a left-only chain, a right-only chain, and a
   balanced tree whose deepest leaves occur on both sides.}
  {I ask each child for its subtree depth and return one plus the larger
   result. Every node is visited once, while extra space equals the tree
   height because of recursion.}

\chaptercheckpoint
  {A parent answer depends on the deeper child subtree.}
  {Recursive subtree depths.}
  {Counting edges when the requested depth counts nodes.}
  {$O(n)$ time and $O(h)$ stack space.}


\problemstart{Kth Smallest Element in a BST}{230}{Medium}{Inorder traversal}{BST and stack}

\subsection{The Problem}
Given a binary search tree and \(k\), return its \(k\)-th smallest value.

\subsection{Use the BST Ordering}
An inorder traversal of a BST visits values in ascending order. An explicit
stack pauses the walk while descending left, then emits one node at a time.
The \(k\)-th popped node is the answer.

\begin{invariantbox}{the next popped node is the smallest unvisited value}
The stack stores the path to the next inorder node. All values already popped
are smaller, and every unvisited value in a right subtree is processed only
after its parent.
\end{invariantbox}

\begin{visualbox}{Inorder turns the BST into sorted output}
\centering
\begin{tikzpicture}[font=\small,>={Stealth[length=2.2mm,width=1.55mm]}]
  \node[trienode] (r) {3}; \node[trienode,below left=11mm and 15mm of r] (a) {1};
  \node[trienode,below right=11mm and 15mm of r] (b) {4};
  \node[trienode,below right=10mm and 6mm of a] (c) {2};
  \draw[flowarrow] (r)--(a); \draw[flowarrow] (r)--(b); \draw[flowarrow] (a)--(c);
  \node[dsnode,right=28mm of r,align=left] {inorder: \(1,2,3,4\)\\\(k=2\Rightarrow2\)};
  \node[draw=orange,very thick,fit=(c),inner sep=1mm] {};
\end{tikzpicture}
\end{visualbox}

\subsection{Pseudocode}
\begin{lstlisting}[style=pseudocode]
stack = empty; current = root
while current exists or stack is not empty:
    while current exists:
        push current; current = current.left
    current = pop stack
    decrement k; if k is zero: return current.value
    current = current.right
\end{lstlisting}

\subsection{C++ Implementation}
\begin{lstlisting}[style=compactcode,caption={Complete solution: Kth Smallest in a BST}]
#include <stack>

using namespace std;

class Solution {
public:
    int kthSmallest(TreeNode* root, int k) {
        stack<TreeNode*> path;
        TreeNode* current = root;
        while (current != nullptr || !path.empty()) {
            while (current != nullptr) {
                path.push(current);
                current = current->left;
            }
            current = path.top();
            path.pop();
            if (--k == 0) return current->val;
            current = current->right;
        }
        return -1;
    }
};
\end{lstlisting}

\subsection{Time and Space}
The traversal stops after reaching the answer. Worst-case time is \(O(n)\),
and stack space is \(O(h)\). The problem guarantees a valid \(k\), so the
fallback return is unreachable for valid input.

\lessonexpansion
  {For the BST \texttt{[3,1,4,null,2]} with \(k=2\), push 3 and 1 while
   descending left. Pop 1 as rank one, move to its right child 2, then pop 2
   as rank two and return it.}
  {Inorder traversal emits BST values in sorted order. The stack stores the
   unprocessed ancestors needed to resume that order, so the \(k\)-th pop is
   exactly the \(k\)-th smallest value.}
  {Check \(k=1\), \(k=n\), a skewed tree, a balanced tree, and a target found
   after entering a right subtree.}
  {I perform iterative inorder traversal and decrement \(k\) whenever a node
   is popped. Since inorder is sorted for a BST, the node that reduces
   \(k\) to zero is the answer.}

\chaptercheckpoint
  {A rank is requested from a binary search tree.}
  {An inorder path stack and a remaining rank counter.}
  {Using preorder, which does not emit sorted values.}
  {$O(n)$ worst-case time and $O(h)$ space.}


\problemstart{Lowest Common Ancestor of a BST}{235}{Medium}{Ordered tree search}{BST}

\subsection{The Problem}
Return the lowest node in a BST that is an ancestor of both given nodes. A
node may be an ancestor of itself.

\subsection{Follow the Shared Direction}
If both target values are smaller than the current value, their ancestor lies
left. If both are larger, it lies right. Otherwise the paths split at the
current node, which is their lowest common ancestor.

\begin{invariantbox}{both targets remain inside the current subtree}
Every move selects the only child subtree that can contain both values under
the BST ordering. The first node where the targets no longer share a strict
direction is the deepest node whose subtree contains them both.
\end{invariantbox}

\begin{visualbox}{The first path split is the answer}
\centering
\begin{tikzpicture}[font=\small,>={Stealth[length=2.2mm,width=1.55mm]}]
  \node[trienode,active] (r) {6};
  \node[trienode,visited,below left=12mm and 20mm of r] (a) {2};
  \node[trienode,visited,below right=12mm and 20mm of r] (b) {8};
  \node[trienode,below left=11mm and 8mm of a] (c) {0};
  \node[trienode,below right=11mm and 8mm of a] (d) {4};
  \draw[flowarrow] (r)--(a); \draw[flowarrow] (r)--(b);
  \draw[flowarrow] (a)--(c); \draw[flowarrow] (a)--(d);
  \node[above=5mm of r,text=orange]{LCA of 2 and 8};
\end{tikzpicture}
\end{visualbox}

\subsection{Pseudocode}
\begin{lstlisting}[style=pseudocode]
current = root
while current exists:
    if p and q are both smaller: current = current.left
    else if p and q are both larger: current = current.right
    else: return current
return null
\end{lstlisting}

\subsection{C++ Implementation}
\begin{lstlisting}[style=compactcode,caption={Complete solution: LCA of a BST}]
using namespace std;

class Solution {
public:
    TreeNode* lowestCommonAncestor(TreeNode* root,
                                   TreeNode* p, TreeNode* q) {
        TreeNode* current = root;
        while (current != nullptr) {
            if (p->val < current->val && q->val < current->val) {
                current = current->left;
            } else if (p->val > current->val && q->val > current->val) {
                current = current->right;
            } else {
                return current;
            }
        }
        return nullptr;
    }
};
\end{lstlisting}

\subsection{Time and Space}
Only one root-to-node path is followed: \(O(h)\) time and \(O(1)\) space.
Unlike the general binary-tree version, no parent map or full traversal is
needed because ordering tells which subtree can contain each value.

\lessonexpansion
  {In the BST rooted at 6, targets 2 and 4 are both smaller than 6, so move to
   node 2. At node 2, one target equals the current value while the other lies
   to the right; their paths split here, making node 2 the LCA.}
  {BST ordering proves that when both targets are smaller or both are larger,
   only one child subtree can still contain both. The first node where neither
   shared direction applies is their deepest common ancestor.}
  {Check targets on opposite sides of the root, one target equal to the LCA,
   adjacent nodes, and both targets deep in the same subtree.}
  {I follow the one direction shared by both target values. When they split,
   or when the current node equals one target, the current node is the lowest
   common ancestor.}

\chaptercheckpoint
  {Find a shared ancestor in an ordered binary tree.}
  {The current root and the two target values.}
  {Ignoring the case where the current node equals one target.}
  {$O(h)$ time and $O(1)$ space.}


\patternintoc{Heaps and Scheduling}
\problemstart{Meeting Rooms II}{253}{Medium}{Minimum end-time heap}{Intervals and min heap}

\subsection{The Problem}
Return the minimum number of rooms needed to hold all meetings without
overlap.

\subsection{Track Rooms by Their End Times}
Sort meetings by start. A min heap stores the ending time of every room
currently in use. Before assigning a meeting, release every room whose meeting
has ended; then push the new end. The maximum heap size is the answer.

\begin{invariantbox}{the heap contains exactly the occupied rooms}
After expired end times are removed for a start time, each remaining heap
entry belongs to a meeting overlapping the new one. Pushing the current end
makes the heap size equal to the rooms simultaneously occupied.
\end{invariantbox}

\begin{visualbox}{The earliest ending room is released first}
\centering
\begin{tikzpicture}[font=\small,>={Stealth[length=2.2mm,width=1.55mm]}]
  \node[dsnode,active] (m) {new meeting\\starts at 5};
  \node[trienode,visited,right=25mm of m] (h0) {4};
  \node[trienode,below left=10mm and 8mm of h0] (h1) {8};
  \node[trienode,below right=10mm and 8mm of h0] (h2) {10};
  \draw[-,thick,navy] (h0)--(h1); \draw[-,thick,navy] (h0)--(h2);
  \draw[flowarrow,teal] (m)--node[edgelabel,above]{\(5\ge4\)} (h0);
  \node[dsnode,fill=softteal,below=8mm of m] (free) {pop 4\\reuse its room};
  \draw[flowarrow,teal] (m)--node[edgelabel,left,font=\scriptsize]{then} (free);
  \node[above=6mm of h0]{min-heap of room end times};
\end{tikzpicture}
\end{visualbox}

\subsection{Pseudocode}
\begin{lstlisting}[style=pseudocode]
sort meetings by start time
minHeap = empty; answer = 0
for each meeting:
    while heap top <= meeting.start: pop heap
    push meeting.end
    answer = max(answer, heap size)
return answer
\end{lstlisting}

\subsection{C++ Implementation}
\begin{lstlisting}[style=compactcode,caption={Complete solution: Meeting Rooms II}]
#include <algorithm>
#include <functional>
#include <queue>
#include <vector>

using namespace std;

class Solution {
public:
    int minMeetingRooms(vector<vector<int>>& intervals) {
        sort(intervals.begin(), intervals.end());
        priority_queue<int, vector<int>, greater<int>> ends;
        int answer = 0;
        for (const auto& meeting : intervals) {
            while (!ends.empty() && ends.top() <= meeting[0]) ends.pop();
            ends.push(meeting[1]);
            answer = max(answer, static_cast<int>(ends.size()));
        }
        return answer;
    }
};
\end{lstlisting}

\subsection{Time and Space}
Sorting and heap operations cost \(O(n\log n)\) time. The heap uses \(O(n)\)
space in the worst case. Removing all expired rooms keeps the key rule exact;
removing only one is enough for some variants but less directly models
the current occupancy.

\lessonexpansion
  {For \texttt{[[0,30],[5,10],[15,20]]}, push 30 for the first room. At start
   5, room 30 is still occupied, so push 10 and the heap size becomes two. At
   start 15, pop 10 because it has ended, then push 20. The maximum heap size
   is two rooms.}
  {After expired end times are removed, every heap entry corresponds to a
   meeting that overlaps the new start. So the heap size is exactly
   the current room demand, and its maximum is the minimum number of rooms
   that can schedule all meetings.}
  {Check no meetings, one meeting, touching meetings, completely nested
   meetings, identical intervals, and many meetings ending at the same time.}
  {I sort by start time and keep active end times in a min heap. The earliest
   ending rooms are released before each insertion. The maximum number of
   active end times is the required room count.}

\chaptercheckpoint
  {Return maximum simultaneous interval occupancy.}
  {Meetings sorted by start and a min heap of active end times.}
  {Using a max heap, which hides the room that becomes free first.}
  {$O(n\log n)$ time and $O(n)$ space.}
